\documentclass[11pt,a4paper]{article}

\usepackage[T1]{fontenc}
\usepackage[utf8]{inputenc}
\usepackage[margin=1in]{geometry}
\usepackage{amsmath,amssymb}
\usepackage{microtype}
\usepackage{xurl}
\usepackage[numbers,sort&compress]{natbib}
\usepackage[hidelinks]{hyperref}

\hypersetup{pdftitle={Serving Image and Video DiTs on Bandwidth-Constrained PCIe Nodes}}
\urlstyle{same}
\setlength{\emergencystretch}{2em}
\newcommand{\evidence}[1]{\footnote{Experimental record in the accompanying
\texttt{dit-rewrite} repository: \nolinkurl{#1}.}}

% References use the companion dit_serving_references.bib database.
% Build from this directory with a LaTeX/BibTeX-capable toolchain.
% Evidence was checked against the completed reports on 2026-10-14; the
% 2026-10-18 additions cite results/a100-e* (local A100 PCIe 40GB). The
% 2026-10-19 revision replaces the 14B extrapolation and the assumed-link
% simulator study with measurements on the target 8xA100 PCIe node, whose
% records are the K-numbered rows of docs/kairos_preexperiments.md.
% Keep the displayed equations: Gaussian noising, a clean-latent estimate,
% and the parallel-form cost model.
% The noise-prediction example does not specify Wan's actual solver update.

\title{Serving Image and Video Diffusion Transformers\\on Bandwidth-Constrained PCIe Nodes\\[0.3em]
\large Background, Motivation, Design, and Evaluation Plan (working draft)}
\author{}
\date{}

\begin{document}
\maketitle

\section{Background}
\label{sec:background}

\subsection{Diffusion Models}
\label{sec:diffusion-models}

Diffusion models generate data by learning to reverse a process that
progressively corrupts clean samples with noise. Early successful
formulations, including denoising diffusion probabilistic models (DDPMs),
predominantly employed convolutional U-Net denoisers~\cite{ho2020ddpm}.
Latent diffusion subsequently reduced the computational burden by
performing the iterative generative process in the compressed representation
of a pretrained autoencoder. A decoder then maps the generated latent back
to image space. This separates perceptual compression from generative
modeling while retaining the repeated denoising
structure~\cite{rombach2022ldm}.

Diffusion Transformers (DiTs) replace the U-Net backbone with a Transformer
operating on patches of the latent representation~\cite{peebles2023dit}.
These patches become tokens processed through repeated attention and
feed-forward blocks, with conditioning information injected into the
network. For video generation, tokens additionally represent temporal
structure, increasing the amount of computation within each network
evaluation. The denoising architecture and the generative formulation are
separate choices: adopting a Transformer does not prescribe a particular
noise schedule, prediction parameterization, or numerical sampler.

The basic diffusion process adds Gaussian noise to a clean latent $x_0$:
\begin{equation}
  x_t=\alpha_t x_0+\sigma_t\epsilon.
  \label{eq:noising}
\end{equation}
Here, $\epsilon\sim\mathcal N(0,I)$, while $\alpha_t$ and $\sigma_t$
control the signal and noise levels. Under noise prediction, the model
estimates $\hat\epsilon_\theta$ from the noisy latent, timestep, and condition.
For $\alpha_t>0$, this gives a clean-latent estimate:
\begin{equation}
  \hat x_0=\frac{x_t-\sigma_t\hat\epsilon_\theta}{\alpha_t}.
  \label{eq:denoising}
\end{equation}
These equations summarize noising and noise prediction; the sampler
determines the iterative updates used during generation~\cite{ho2020ddpm}.

For text-conditioned latent generation, inference begins by encoding the
prompt and initializing a Gaussian latent. Each sampling iteration
evaluates the Transformer at the current latent and timestep, combines
conditional predictions when required, and updates the latent through the
sampler. Classifier-free guidance (CFG) combines conditional and
unconditional predictions using a guidance scale~\cite{ho2022cfg}; its
two branches introduce additional network evaluations that can be batched
or otherwise optimized.

The Wan model used in our experiments employs a flow-matching formulation
that predicts a velocity field along a path between Gaussian noise and
data~\cite{wan2025}. Our implementation uses UniPC to update the latent
from velocity predictions, rather than applying the noise-prediction
example above. The serving abstraction remains repeated \emph{model
prediction followed by sampler update}, including any solver history.
After the final update, the VAE decoder produces the output video. This
formulation exposes the execution hierarchy relevant to serving: a request
contains sampling iterations, an iteration may contain multiple guidance
branches, and each branch contains a sequence of Transformer blocks.

\subsection{Diffusion Serving}
\label{sec:diffusion-serving}

A diffusion serving system executes requests arriving over time while
sharing limited compute and memory resources. A request specifies its
conditioning input, model, output dimensions, sampling configuration, and
potentially a completion deadline. Its end-to-end latency extends from
arrival to output completion, encompassing encoding, queueing, execution,
interruptions, and decoding. This differs from isolated generation time
measured after inputs and models have already been prepared.

Throughput measures completed requests per unit time, whereas
deadline-constrained goodput measures successfully completed requests that
meet their deadlines per unit time. End-to-end latency is the difference
between a request's completion and arrival times. For finite traces, our
measurement interval includes idle periods and the final drain. Failed or
late requests do not contribute to goodput, while their execution time
remains in the observation interval. Output quality is evaluated
separately under the prescribed comparison protocol.

Three properties shape the scheduling problem. First, image and video
requests can differ substantially in service time, even when both use
Transformer-based generators. Second, repeated sampling creates
opportunities to suspend execution between iterations, guidance branches,
or smaller network segments. Third, these boundaries expose different
amounts of live state. A step boundary may require the current latent and
solver history; an intra-step boundary can additionally require
intermediate activations, branch outputs, and execution position.

Feature caching introduces another category of state whose lifetime spans
sampling iterations. Cached features can reduce subsequent computation,
but preserving, transferring, or discarding them changes memory occupancy
and future execution cost. Consequently, releasing the GPU for another
request does not necessarily release the suspended request's memory.

Batching, preemption, and parallel execution address different aspects of
this problem. Batching amortizes execution costs across compatible
requests; preemption limits how long one request blocks another;
parallelism distributes computation across devices. Their effectiveness
depends on workload composition, execution granularity, and state
requirements. A serving policy must therefore evaluate both immediate
responsiveness and the cost of completing suspended work.

\section{Motivation}
\label{sec:motivation}

Our measurements investigate mixed image--video serving on an A100 PCIe
GPU with 40\,GB of memory. The principal configuration combines
Wan2.1-T2V-1.3B video generation with Sana-Sprint-0.6B image generation.
Wan uses 32 sampling steps, CFG scale 5, and Flash attention; its text
encoder runs on the CPU, while its DiT and VAE reside on the GPU. We
distinguish local execution measurements from online experiments
throughout.

\subsection{Preemption Granularity}
\label{sec:temporal-heterogeneity}

A sampling iteration can already be too coarse to protect short requests.
At $480\times832$ resolution, the measured average Wan step time increases
from 0.581 seconds for 17 frames to 2.021 seconds for 49 frames and 4.122
seconds for 81 frames. Each estimate comprises six measured steps from
truncated profiling runs. By comparison, Sana generates a complete
$1024\times1024$ image in approximately 0.317 seconds under its two-step
configuration. Thus, one video sampling step can occupy the GPU for
roughly thirteen times the service time of a complete short image
request.\evidence{results/research-directions/run01/direction1/direction1-profile.json}

This mismatch limits step-boundary preemption. An urgent image request
arriving shortly after a video step begins must wait for the remaining
step even if the scheduler immediately recognizes its priority. We
therefore evaluate suspension at CFG boundaries and after individual
Transformer blocks or groups of blocks. In a paired microbenchmark
comprising two rounds and three arrival positions per round, the
descriptive p95 waiting time falls from 3,732\,ms with step-boundary
suspension to 1,659\,ms at CFG boundaries, 44\,ms after individual blocks,
and 168\,ms after groups of four blocks. These measurements use preencoded
image conditions and a three-step video prefix, with six observations per
policy.\evidence{results/research-directions/run01/direction1/direction1-preemption.json}

Finer suspension also retains additional state. For the 81-frame
configuration, block-level suspension requires approximately 117.5\,MiB
of extra intermediate state. Separate correctness experiments interleave
image execution with paused video inference and verify the resumed
trajectory against uninterrupted execution. All eight tested cases pass,
covering two video shapes. This establishes the feasibility of the tested
execution mechanism; the scheduling benefit still requires workload-level
evaluation.\evidence{results/research-directions/run01/direction1/direction1-correctness.json}

The available online results support further investigation. Eleven traces
have complete comparisons across all five policies, giving 264 requests
per policy: 242 images and 22 videos. Step-boundary scheduling completes
111 requests on time, compared with 202 for block-1 and 199 for block-4;
all 22 videos meet their deadlines in these three arms. However, the full
experiment stopped after 57 of 60 planned replays because its allocated
budget was exhausted. These are descriptive results from the complete
paired subset. They motivate selecting suspension granularity according
to short-request response requirements, while accounting for retained
state and switching
overhead.\evidence{results/research-directions/run01/REPORT.md}

\subsection{Cache State}
\label{sec:cache-state}

Feature caching changes the computational structure a scheduler sees and
trades output for time. A CFG-coordinated cache cut the service time of an
81-frame Wan2.1-1.3B request by 19.6\% within a mean-LPIPS screen of 0.05,
with reuse steps about $3.75\times$ cheaper than recompute
steps.\evidence{results/research-directions/run01/direction2/qualify.json}
The cached state of a suspended request is then a scheduling object of its
own: offloading 448\,MB to the CPU takes 197\,ms and restores in 38\,ms
with identical output, whereas clearing it is free but raises the next
step from 1.13 to 4.19\,s and changes the output; keeping the state of up
to 16 paused requests never blocked an urgent one on a 40\,GB
GPU.\evidence{results/research-directions/run01/direction2/state_cost.json}
We leave caching out of the serving design that follows: block quanta
require every block to be computed, cached features do not survive a
change of parallel form, and every cache configuration would need its own
quality qualification. The measurements remain as the cost of that choice
and as future work on cache-aware suspension.

\subsection{Serving Performance}
\label{sec:local-versus-serving}

Improving the execution time of one model is valuable, but the serving
objective depends on how that improvement interacts with arrivals,
deadlines, and competing requests. We examine this distinction by
compiling Wan Transformer blocks while retaining the existing block-4
scheduling boundaries. Across four distinct inputs measured in two
execution orders, compilation reduces local service time by
14.40--14.44\%. These measurements include preparation after preencoding,
denoising, and decoding, with eight observations per arm.

The corresponding online experiment uses two frozen arrival traces and
both execution orders, producing eight complete replays and 192 request
executions. Each arm contains 96 requests, including 88 images and eight
videos. Despite the local reduction, pooled goodput decreases from
0.03773 to 0.03709 requests per second, a 1.70\% decline. The pooled
image-latency p95 increases from 1.430 to 1.627 seconds, or 13.76\%. All
eight videos in each arm meet their deadlines. These finite measurements
demonstrate that the observed local acceleration does not translate into
improved serving outcomes under the tested
configuration.\evidence{results/optimization-screen/run01/wan_compile/attempt03/online.json}

The aggregation is important. Pooled goodput divides the total number of
on-time completions by the sum of replay durations; it does not average
per-replay goodput values. Similarly, the image p95 is recomputed from
the 88 image latencies in each arm. These observations do not establish
statistically significant degradation or identify its dominant cause.
They establish the need to evaluate execution optimizations against the
actual online objective using unchanged arrivals, deadlines, and complete
accounting.

Together, the measurements motivate a controller that considers request
urgency, legal execution boundaries, and continuation costs. Fine-grained
earliest-deadline-first scheduling provides a strong baseline: additional
rules must demonstrate benefits beyond those obtained simply by exposing
shorter execution intervals. The following measurements extend this
observation to the multi-GPU setting that motivates our design.

\subsection{Long Steps and Multi-GPU Execution}
\label{sec:long-steps}

Step boundaries become rarer as models and outputs grow. On the same A100,
one Wan2.1-1.3B solver step takes 0.58\,s at $480\times832$ with 17 frames,
4.16\,s with 81 frames, and 16.8\,s at $720\times1280$ with 81 frames; each
step consists of 60 Transformer-block evaluations of almost equal cost
(two CFG branches of 30 blocks).\evidence{results/a100-e2-cost-profile/run01}
On the target machine, an $8\times$A100 PCIe 40\,GB node without NVLink
(every GPU pair crosses a host bridge), Wan2.2-TI2V-5B takes 8.15\,s per
step at its official $704\times1280$, 121-frame output and 1.93\,s at
$512\times896$ with 81 frames, and its FP32 VAE decode takes 32.9\,s and
10.3\,s respectively; a Sana-Sprint image takes
0.32\,s.\evidence{docs/kairos\_preexperiments.md (K17)}
A 50-step 720p request therefore occupies one GPU for about 440\,s, during
which roughly 1,400 image requests could have been served.
Memory is the second constraint of a 40\,GB card: the FP32 decode of the
720p request peaks at 32\,GB inside the video worker, which cannot coexist
with the 8\,GB image worker on one GPU even when both return their
allocator caches after every request; decoding in BF16 halves the peak to
21\,GB and the time to 17\,s, and changes the decoded video by a mean LPIPS
of 0.00025 (57\,dB) relative to FP32 from identical latents, so the 5B
serving experiments decode in
BF16.\evidence{docs/kairos\_preexperiments.md (K20, K21)}
Sequence parallelism shortens a step, but the request then holds several
GPUs, so a step-boundary scheduler must either wait for the whole group or
keep GPUs idle for short requests.

The effect survives the full serving path on one GPU. Through the native
multi-GPU executor with real Wan and Sana workers, a 30-request Poisson mix
of images and 17/33-frame videos at loads of 0.4, 0.6, and 0.8 of measured
capacity meets 73\%, 67\%, and 50\% of deadlines under step-boundary EDF,
and 93\%, 80\%, and 83\% when videos advance in four-block
quanta.\evidence{results/a100-e3-load-curves/run01}
Two costs of that design are visible in the same runs. Feasible-first
ordering lets requests that are already late wait behind feasible videos,
raising the image p95 latency to 74--118\,s at high load; keeping late
requests with at most 2\,s of remaining work in EDF order removes this tail
(image p95 1.8\,s) and raises on-time completions under block quanta from 24
to 27 of 30.\evidence{results/a100-e9-late-short/run01}
The lazy variant, which lets an image wait for a nearby step boundary,
gained nothing over eager preemption on one GPU.

\paragraph{Sequence parallelism on a PCIe node.}
Ulysses sequence parallelism (SP) exchanges the query, key, value, and
output projections of every attention layer through all-to-all collectives.
On the target node these collectives run at 2.3--4.8\,GB/s per rank, and
the host bridge caps the whole node at about 19\,GB/s: when four or more
GPUs communicate at once each gets 2.4\,GB/s whether or not other groups are
active, and two concurrent pairs fall from 4.8 to
3.9\,GB/s.\evidence{docs/kairos\_preexperiments.md (K3, K15)}
Measured whole-step speedups are correspondingly small. For Wan2.2-5B at
$704\times1280\times121$, SP2 gives $1.36\times$, SP4 $1.7$--$1.8\times$, and
SP8 $3.1\times$; at $512\times896\times81$ the figures are $1.13$, $1.29$, and
$2.09\times$; for Wan2.1-1.3B at 720p with 81 frames SP4 gives
$2.09\times$, and 17-frame clips gain nothing at any
degree.\evidence{docs/kairos\_preexperiments.md (K7, K14, K15)}
Communication therefore accounts for 32--75\% of an SP step across the
shapes a service would see. A three-term model,
$t_{\mathrm{step}}(d)=t_{\mathrm{compute}}/d+V\,(d-1)/(d^2 B_d)+c\,\ell_d$,
with $V$ the all-to-all volume of one step, $B_d$ and $\ell_d$ the per-rank
rate and latency of a $d$-GPU group, and $c$ the number of collectives,
fitted on the 1.3B model ($B_2=4.8$, $B_4=B_8=2.7$\,GB/s) predicts every
5B shape and degree within 3.7\% (SP8 within 2.6\%) without
refitting.\evidence{docs/kairos\_preexperiments.md (K14, K15)}
The rates, however, belong to a placement rather than to the node: two jobs
that received different GPU quartets of the same node measured SP4 steps
of 3.38 and 4.02\,s for the same request, and effective pair rates between
5.1 and 8.7\,GB/s, while \texttt{nvidia-smi} reports every pair as the same
host-bridge class. In-model rates also exceed those of a raw collectives
probe, because the attention backend issues the three input all-to-alls
asynchronously. A serving system therefore has to calibrate its link table
per placement by timing one step of each parallel form on the groups it
owns, which takes about a minute, and price a form that spans several
groups by its slowest
group.\evidence{docs/kairos\_preexperiments.md (K16b, K16d)}

\paragraph{Where the time can go instead.}
The two CFG branches of a step are independent until their predictions are
combined. Placing them on the two halves of a group (\emph{CFG parallel};
xDiT uses the same split statically) replaces the all-to-all traffic of
SP2 (about 10\,GB per rank per step for the 720p request) with one 10.5\,MB
exchange of predictions, and the combination and solver update then run
identically on both halves. On the target node this form is bit-exact with
single-GPU execution and reaches $2.00\times$ on two GPUs and, with SP2
inside each half, $2.56$--$2.70\times$ on four, against $1.33$--$1.55\times$
and $1.71$--$2.44\times$ for plain SP on the same
GPUs.\evidence{docs/kairos\_preexperiments.md (K16, K16b, K16c)}
Within an SP group, the two branches can also be interleaved block by block
on two CUDA streams so that one branch's all-to-all overlaps the other's
compute. The interleaved execution hides 33--87\% of the communication of a
step depending on degree and placement (SP4 reaches $2.4$--$3.0\times$) and
is bit-exact with the plain SP run of the same layout; what plain SP
itself changes relative to one GPU is the token-shard size its GEMMs see,
not anything the interleaving
adds.\evidence{docs/kairos\_preexperiments.md (K16b, K16c, K16e)}
Plain SP is not exact on GPUs: at 50 steps the 5B model's SP4 output at
$512\times896\times81$ differs from the single-GPU output by a mean LPIPS of
0.107 (SP2: 0.041), while at $704\times1280\times121$ both stay at
0.0002.\evidence{docs/kairos\_preexperiments.md (K18)}
CFG forms are therefore the only parallel forms that preserve the output,
and on this hardware they are also the fastest per GPU.

\paragraph{What elasticity buys.}
With real Wan and Sana workers on four target GPUs, block-level suspension
with step-boundary regrouping met 100 of 100 deadlines where static SP1 met
94, static SP2 61, static SP4 32, and step-boundary elastic SP
68.\evidence{docs/kairos\_preexperiments.md (K8)}
At twice the capacity, 200 requests per arm, block quanta kept 153 of 154
images on time against 68 under step quanta with the same planner, while
video completions were equal (36 and 36); elastic SP completed one more
video than SP1-only execution.\evidence{docs/kairos\_preexperiments.md (K11$'$)}
The calibrated simulator reproduces these counts within about three videos
and shows the same pattern over ten seeds: on this hardware elastic SP
raises video attainment by at most 0.003 and shortens the 720p tail (p95
2458 against 2688\,s), whereas block-level preemption is what protects
images.\evidence{docs/kairos\_preexperiments.md (K9b, K10b)}
Choosing among parallel forms by the cost model extends the tail benefit
to the 5B workload: over ten seeds on eight simulated GPUs at the capacity
load, form selection cuts the video p95 by 18\% relative to SP-only
elasticity with every image on time, and under 1.5--2$\times$ overload all
policies are throughput-bound within 0.86--0.90 video attainment, the
static CFG-pair policy among them at the cost of 18\% of the
images.\evidence{docs/kairos\_preexperiments.md (K19)}

Suspension inside a step is cheap in state. With caching disabled, the
step-boundary request state of the measured shapes occupies 16--86\,MB and
restores in under 10\,ms. A block-boundary continuation additionally keeps
28--260\,MiB of live activations.\evidence{results/a100-e2-cost-profile/run01}
Contrary to the view that intermediate points of a step are not valid
stopping points, the continuation can be serialized and resumed in another
process. For all six measured shapes, a Wan2.1-1.3B request suspended after
13 blocks of the conditional branch or 15 blocks of the unconditional branch
(which also carries the conditional prediction) was exported with its
untouched solver state, released, and resumed in a fresh worker process; the
decoded videos and final latents matched an uninterrupted run bit for bit.
The packets occupied 28--269\,MiB; export and import took 0.15--1.3\,s and
0.09--0.45\,s including integrity hashing.\evidence{results/a100-e5-block-exact/run01}
Finally, the VAE decode becomes the remaining non-preemptible tail once
denoising is preemptible: it takes 1.5--15.9\,s for the measured shapes,
but decodes causally one latent frame at a time and can therefore yield
between frames. Decoding in such chunks matches the whole decode bit for bit
and, on a frozen burst trace, raised on-time completions from 7 to 10 of 24
under block quanta and from 5 to 8 under step quanta.\evidence{results/a100-e7-chunked-decode/run01}

Time-sliced co-location is not a substitute for preemption. With
Mistral-7B, Wan2.1-1.3B, and Sana-Sprint resident as three processes on
one 40\,GB A100, a concurrent Wan step slows image generation from 0.30 to
0.58\,s and LLM decoding from 34--41 to 63--66\,ms per token, while the
Wan step itself slows by 6--22\%. The three resident models and one video
leave 7.4\,GiB free; six 3.9k-token conversations (488\,MiB of KV each)
reduce that to 0.53\,GiB, and per-process allocator caches hold about
11\,GiB that the other processes cannot use.\evidence{results/a100-e8-colocation/run01}
On 40\,GB GPUs, placement of suspended DiT state and LLM KV therefore
becomes a real constraint once models share a device.

The present evidence therefore supports a multi-GPU design for PCIe servers
in which requests advance in block quanta, parallel forms are chosen per
request from a placement-calibrated cost model and change only at step
boundaries, communication is hidden or avoided rather than accelerated, and
short requests preempt a long request inside a step only when they cannot
wait for its boundary.

\section{Design Outline}
\label{sec:design}

The system serves image and video requests on one PCIe node. A C++
controller plans every quantum with a cost model; Python workers, one video
and one image worker per GPU, execute quanta and keep request state. The
same planner drives a discrete-event simulator, so a simulated policy and a
real one differ only in what executes the quanta.

\paragraph{Execution hierarchy and preemption points.}
A video request is a solver trajectory of $S$ steps; each step evaluates
$2L$ Transformer blocks (two CFG branches of $L$ blocks), then combines the
two predictions and advances the solver; the request ends with a VAE
decode. The worker exposes three kinds of stopping points. A
\emph{step boundary} holds only the replicated solver state (16--92\,MB for
the measured shapes). A \emph{block boundary} inside a step additionally
holds the running packet of the branch in flight; it is reached by
executing the step as a sequence of quanta of $b$ blocks, and the
continuation can be exported to another process and resumed with
bit-identical output. A \emph{decode chunk} boundary falls between latent
frames of the causal VAE decode. Quanta are time-targeted: the planner
converts a quantum length (300\,ms in our experiments) into a block count
for the request's class and parallel form, so a preemption check happens
about every 300\,ms of GPU work regardless of model or degree.

\paragraph{Parallel forms and the cost model.}
A video quantum runs on an aligned group of $g=c\cdot d$ GPUs in a
\emph{form} $(c,d)$: the $c\in\{1,2\}$ CFG branches run on separate halves
of the group, and sequence parallelism of degree $d$ runs inside each half.
The forms used here are SP1, SP2, SP4, SP8, CFG2, CFG2$\times$SP2, and
CFG2$\times$SP4. Under CFG parallelism each half computes one branch for
the whole step on a replicated state, the halves exchange their predictions
once, and both apply the same combination and solver update, so the state
stays replicated and the output equals single-GPU execution. Under SP the
two branches may be \emph{interleaved}: block $i$ of the conditional branch
runs on the main CUDA stream while block $i$ of the unconditional branch
runs on a second stream, so one branch's all-to-all overlaps the other's
compute; both branches share the rotary and timestep embeddings and stop at
the same block, so quanta and exports are unchanged. The planner prices a
form by
\[
t_{\mathrm{step}}(c,d,n)=\frac{t_{\mathrm{compute}}}{c\,d}
+\frac{V}{c}\cdot\frac{d-1}{d^2 B_d(n)}\,(1-h_d\,[c{=}1])
+\frac{k}{c}\,\ell_d+[c{=}2]\Big(\frac{E}{B_2}+\ell_2\Big),
\]
where $V$ is the all-to-all volume and $k$ the number of collectives of a
step, $E$ the size of one prediction, $B_d(n)$ the per-rank rate of a
$d$-GPU group when $n$ SP groups share the host bridge, $\ell_d$ its
latency, and $h_d$ the fraction of the shorter of the compute and
communication terms that interleaving hides. A fixed per-quantum overhead
(12\,ms measured end to end) is added for every quantum a request still
needs. $B_d(n)$, $\ell_d$, and $h_d$ are calibrated per placement by timing
one step of each form at start-up; everything else is a class constant
derived from the model architecture and one single-GPU profile.

\paragraph{Planner.}
Between quanta no request owns a GPU, and the planner recomputes the
assignment of waiting and suspended requests to free GPUs. Requests are
ordered feasible-first by deadline: a request that would miss its deadline
even on the widest group free at planning time yields to those that can
still make it, except that requests with less than 2\,s of remaining work
keep their EDF position, because demoting them only turns a small miss into
a long tail. A video takes the smallest group whose cheapest form meets its
deadline, or the widest free group when nobody else waits; with other
requests waiting it may still widen to a form whose per-GPU efficiency is
at least a threshold (0.95 here), which admits CFG pairs but not PCIe SP,
because such a form returns its GPUs as soon as it took them. Within a
group size the form with the smallest predicted step time wins. Groups and
forms change only at step boundaries; a request inside a step resumes on its
group, and the planner never advances another request onto a GPU that holds
a mid-step continuation, a contract checked after every plan in both the
executor and the simulator. An image takes an idle GPU if one exists, else
the GPU of a video at a step boundary, and preempts a mid-step group only if
waiting for that group's boundary would miss its deadline (lazy
preemption). Finished videos decode on the first GPU of their group in
latent-frame chunks of the same time target.

\paragraph{State movement.}
A video prepared on one GPU joins a wider group at a step boundary through
a host-staged export of its replicated state (about 16\,MiB per move for
the 1.3B model, 70--150\,ms), after which every rank of the new group
holds the same state; a CFG pair needs no further transfer because both
halves keep the full state. Prompt encoding runs on the CPU ahead of
admission and is never on the GPU path.

\paragraph{What is deliberately left out.}
Feature caching is off throughout: block quanta require every block to be
computed, cached features do not survive a change of group, and caching
trades output for time, which the SLO evaluation would then have to qualify
per configuration. Image batching is not used by any policy. Requests never
span nodes.

\section{Evaluation Plan}
\label{sec:evaluation-plan}

\emph{One $8\times$A100 PCIe 40\,GB node (no NVLink), the hardware of
Section~\ref{sec:long-steps}.}
Models: Wan2.2-TI2V-5B at $704\times1280\times121$ (``long'', 50 steps) and
$512\times896\times81$ (``medium''), the model GENSERVE and GF-DiT also
use; Wan2.1-1.3B as a short-step video model with a fully calibrated
simulator; Sana-Sprint-0.6B for images. Workloads: Poisson arrivals at
1$\times$ and 1.5$\times$ the single-GPU capacity of the node, video:image
mixes of 20/80 and 50/50 by request count, deadlines of $1.5\sigma$ times
the standalone single-GPU service time for videos with
$\sigma\in\{0.8,1.0,1.3\}$ and $2.5\times$ for images, plus a burst trace
of images arriving during long videos. Metrics: per-class deadline
attainment, goodput, p50/p95/p99 latency, and GPU-hours; one seed per real
cell, a second seed on the headline cells, ten seeds in the simulator.

\paragraph{Baselines.}
All arms share the planner, workers, quanta, and cost model, so they differ
in one rule each: static SP1/SP2/SP4, static CFG2 (xDiT-like fixed
hybrid), step-boundary elastic SP (the GENSERVE/GF-DiT regime, without
batching), round-based regrouping every $r$ steps (TetriServe-like), and
EDF/SRTF/FCFS ordering.

\paragraph{Ladder.}
Each contribution is isolated by adding one mechanism at a time to the
step-boundary SP-only baseline, on the real node (one seed) and in the
simulator (ten seeds): (1) block quanta with lazy preemption, (2) chunked
decode, (3) interleaved branches, (4) form selection with CFG forms,
(5) contention-aware rates. Component ablations on the full system turn off
one item each: lazy versus eager preemption, free versus widest lateness
test, the late-short rule, quantum length (150/300/600\,ms), an image
reservation, and contention-blind planning.

\paragraph{Simulator validation.}
The simulator is calibrated once from the 1.3B cluster runs (contended
pair rate and per-quantum overhead) and then checked out of sample: on the
4-GPU 2$\times$-load trace it predicted the on-time counts of four
policies in the measured order, each about three videos too high, with
per-step times within 2.5\%; on the 5B node it is checked the same way on
the 4-GPU replay before the sweeps.\evidence{docs/kairos\_preexperiments.md (K9b, K11$'$)}
Sweeps beyond the real cells (loads up to 2$\times$, mixes to 80/20,
$\sigma$ from 0.8 to 1.3, and link rates from 4.8 to 20\,GB/s to locate
where elastic SP starts to pay) are reported with that bias.

\section{Related Work}
\label{sec:related-work}

\subsection{Diffusion Serving Systems}
\label{sec:related-serving}

TetriServe combines deadline-aware scheduling with elastic sequence
parallelism and supports cross-request latent reuse through
NIRVANA~\cite{lu2026tetriserve}.
DiffServe uses query-aware model cascades and resource allocation to
balance generation quality and serving efficiency~\cite{ahmad2025diffserve}.
NIRVANA reuses intermediate noise states across similar prompts to skip
early sampling steps~\cite{agarwal2024nirvana}.

DistriFusion uses displaced patch parallelism and features from previous
diffusion steps to overlap communication with
computation~\cite{li2024distrifusion}.
PipeFusion combines patch and layer partitioning with cross-step feature
reuse to reduce communication in pipelined DiT
inference~\cite{fang2025pipefusion}.
GENSERVE, a preprint, combines step-boundary preemption, elastic sequence
parallelism, dynamic batching, and SLO-aware scheduling for heterogeneous
diffusion workloads~\cite{ye2026genserve}.
FlashDiff, also a preprint, schedules regional computation at patch
granularity while restricting preemption to denoising-step
boundaries~\cite{qiao2026flashdiff}.
GF-DiT schedules per step with group-free collectives so that SP groups can
be re-formed quickly, and argues that points inside a step are not valid
stopping points~\cite{gfdit2026}; our block-boundary continuations are
exported and resumed bit-exactly across processes and SP groups on two
models, which contradicts that premise. GENSERVE evaluates on the same
Wan2.2-5B model and the same hardware class (a PCIe node without NVLink)
but preempts only at step boundaries, chooses SP degrees from a profile
table, and does not model link contention or consider CFG parallelism.
xDiT provides the parallel inference strategies, including the static CFG
split and the hybrid of CFG and sequence parallelism, and already notes
that Ulysses all-to-all is sensitive to PCIe topology~\cite{fang2024xdit};
a CFG-and-pipeline hybrid for two PCIe GPUs reports $2.1$--$2.3\times$ on
image models for a single request~\cite{cfgpipeline2026}. We take the CFG
split into the serving loop as one form among several, choose it per
request from a placement-calibrated cost model, and hide SP communication
losslessly by interleaving the branches rather than by reusing stale
activations as DistriFusion and PipeFusion do. PixelFlow partitions
image requests at token granularity for throughput~\cite{pixelflow2026}.
In LLM serving, ConServe and FlowPrefill preempt at layer or operator
granularity within an iteration~\cite{qiu2024conserve,flowprefill2026}.

\subsection{GPU Sharing and Caching}
\label{sec:related-state}

REEF exploits DNN kernel idempotence for reset-based preemption and uses
dynamic kernel padding to improve GPU sharing~\cite{han2022reef}.
Nixie coordinates kernel dispatch and memory migration for transparent
application-level multiplexing, waiting for outstanding kernels to
finish before moving memory~\cite{xu2026nixie}.

Pyramid Attention Broadcast reuses attention outputs across sampling
steps with intervals tailored to different attention
types~\cite{zhao2025pab}.
FasterCache combines dynamic feature reuse with CFG-Cache to exploit
redundancy between conditional and unconditional
outputs~\cite{lyu2025fastercache}.
TeaCache uses timestep-modulated inputs to estimate output variation and
select reuse opportunities~\cite{liu2025teacache}.
TaylorSeer predicts future features through Taylor expansions using
derivatives estimated from feature history~\cite{liu2025taylorseer}.

These techniques motivate examining how suspension boundaries and cached
state jointly affect responsiveness, memory residency, and continuation
cost; our measurements characterize these tradeoffs without establishing
general superiority over existing systems.

\bibliographystyle{unsrtnat}
\bibliography{dit_serving_references}
\end{document}
